\documentclass{article}

\usepackage[utf8]{inputenc}
\usepackage[german]{babel}
\usepackage{listings}
\usepackage{xcolor}

\lstset{
    basicstyle=\ttfamily\small,
    breaklines=true,
    backgroundcolor=\color{gray!10}
}

\begin{document}
    \begin{center}
        \Large \textbf{Rabbit Challenge 1: User Management} \\ [6pt]
        \normalsize
        Autor: Maximilian Jakob Maag B.Sc. \\
        Version: 1.0
    \end{center}

    \section{Intro}
    In dieser Challenge geht es um die Verwaltung von Benutzern auf einem Linux System mithilfe von Ansible. Ziel ist es, typische Administrationsaufgaben zu automatisieren.

    \section{Aufgaben}
    
    \subsection{Aufgabe 1}
    Erstellen Sie ein Ansible-Playbook, das 5 Benutzer aus einer Variablendatei anlegt. Jeder Benutzer soll ein eigenes Home-Verzeichnis und eine Standard-Shell (/bin/bash) erhalten.
    
    \textbf{Beispiel Playbook:}
    
    \begin{lstlisting}[language=yaml]
- name: Benutzer erstellen
  hosts: all
  become: true

  vars_files:
    - users.yml

  tasks:
    - name: Erstelle Benutzer
      ansible.builtin.user:
        name: "{{ item.name }}"
        shell: /bin/bash
        create_home: yes
      loop: "{{ users }}"
    \end{lstlisting}

    \textbf{Beispiel Variablendatei (users.yml):}
    
    \begin{lstlisting}[language=yaml]
users:
  - name: user1
  - name: user2
  - name: user3
  - name: user4
  - name: user5
    \end{lstlisting}

    \subsection{Aufgabe 2}
    Setzen Sie eine Passwortablaufrichtlinie für alle Benutzer (z.B. 90 Tage Gültigkeit, 7 Tage Warnung).

    \subsection{Aufgabe 3}
    Fügen Sie Benutzer automatisch bestimmten Gruppen zu (z.B. developers, admins).

    \subsection{Aufgabe 4}
    Konfigurieren Sie sudo-Rechte so, dass nur die Gruppe „admins“ Zugriff erhält.

    \subsection{Aufgabe 5}
    Sperren Sie Benutzerkonten, die sich seit 30 Tagen nicht angemeldet haben.

    \subsection{Aufgabe 6}
    Löschen Sie einen Benutzer inklusive Home-Verzeichnis vollständig.

    \subsection{Aufgabe 7}
    Hinterlegen Sie SSH Public Keys für mehrere Benutzer.

    \subsection{Aufgabe 8}
    Erstellen Sie einen Report aller vorhandenen Benutzer (UID, Gruppen, Shell).

    \subsection{Aufgabe 9}
    Stellen Sie sicher, dass alle Benutzer /bin/bash als Standard-Shell verwenden.

    \subsection{Aufgabe 10}
    Legen Sie Benutzer mit einem Ablaufdatum an (z.B. für temporäre Accounts).

    \subsection{Aufgabe 11}
    Deaktivieren Sie den Root-Login über SSH und erlauben Sie nur Key-basierte Authentifizierung.

\end{document}